\documentclass[class=NCU_thesis, crop=false]{standalone}
\begin{document}

\chapter{實驗結果與討論}

\section{相位調製對光強的影響}
我們先以兩種方式去確認相位調製對光強的影響，第一種是用功率計 (power meter) ，去分別測量兩台 EOM 都開啟與關閉時的光強，再將兩個數值相除得到變化率；另一種方法是透過 Fabry-Perot 測量頻譜，比較相位調製前後的總面積大小，測量結果如下。

\fig[0.75][fig:label][!htb]{eom-transmission.png}[EOM 開啟前後之功率變化率]

從隨機訊號 1 Gb/s 到 10 Gb/s 的調製，看不出有特定的變化趨勢，可見展頻與壓縮只會影響頻率的分佈，不會明顯改變光強度。

\section{古典光量測}
\subsection{展頻與壓縮}
此實驗部分的實驗使用（光路圖）的架設，若只開啟第一台 EOM，在 10 Gb/s 隨機訊號的調製下可將窄頻雷射光的頻譜展至 10 GHz 寬，但由於我們的使用的 Fabry-Perot FSR 僅 10 GHz，無法涵蓋完整的頻率區間，要想掃出完整展開的頻譜需使用 FSR 20 GHz 以上的才能看到，所以下面會先用 2 Gb/s 的訊號來測試，展頻的結果是否符合理論模擬。

\subsubsection{入射光之頻譜}
在兩台 EOM 都關閉的情況下，可以測到波長 795 nm 窄頻雷射的頻譜，結果如下圖，以此 Fabry-Perot 的解析度掃出的雷射頻寬約為 30 MHz。

\fig[0.5][fig:label][!htb]{temp.png}[窄頻雷射頻譜]

\subsubsection{5 Gb/s 隨機訊號之相位調製}
先以 5 Gb/s 隨機訊號進行相位調製，只開啟第一台能將頻譜展至 $\pm$5GHz，如下圖。

\fig[0.5][fig:label][!htb]{temp.png}[5 Gb/s 訊號之展頻頻譜]

頻譜的形狀大致上與理論相符，但在 $\pm$5GHz 的位置有一個突起的訊號，這是由於隨機訊號的上升與下降時間不夠快所致，若在數值模擬中把隨機訊號加上約 30 ps 的上升與下降時間（如圖），則會出現類似的結果，如圖：

\begin{figure}[!hbt]
    %\captionsetup[subfigure]{labelformat=empty} % 完全隱藏圖號
    \centering
    \subcaptionbox
        {caption\_1
        \label{fig:subfig_fig1}}
        {\includegraphics[width=0.3\linewidth]{temp.png}}
    ~~~~
    \subcaptionbox
        {caption\_2
        \label{fig:subfig_fig2}}
        {\includegraphics[width=0.3\linewidth]{temp.png}}
\end{figure}

% \fig[0.5][fig:label][!htb]{temp.png}[修正前後之隨機訊號]
% \fig[0.5][fig:label][!htb]{temp.png}[修正後展頻頻譜]

此外，還可看出該頻譜的包絡線有週期振盪的訊號，其原因為我們使用的隨機訊號實際上是個重複出現的週期訊號，每個週期有 $2^{31}-1$ 個位元，若把單位週期的位元數調為 $2^{15}-1$ 做可看到週期更小的震盪週期，如圖：

\begin{figure}[!hbt]
    %\captionsetup[subfigure]{labelformat=empty} % 完全隱藏圖號
    \centering
    \subcaptionbox
        {caption\_1
        \label{fig:subfig_fig1}}
        {\includegraphics[width=0.3\linewidth]{temp.png}}
    ~~~~
    \subcaptionbox
        {caption\_2
        \label{fig:subfig_fig2}}
        {\includegraphics[width=0.3\linewidth]{temp.png}}
\end{figure}

當兩台 EOM 同時開啟時，理論上要能將展寬的頻譜還原成調製前的狀態，但從（圖）的實驗結果可以看出，壓縮回來的頻寬較
寬，若只計算中心頻率附近 500 MHz 的頻率區間，與調製前的相比，光強僅約 70\%，造成壓縮效果不佳的可能原因為兩個隨機訊號的形狀不同，上下也不夠對稱，導致無法將相位做反向的調製，使訊號完美還原成最初的狀態。

\paragraph{10 Gb/s 隨機訊號之相位調製}

從 5 Gb/s 的展頻可看出，展寬的頻率區間如理論計算的結果
頻譜還原的訊號，與調製前的頻譜比較如圖，若將雷射調至會被銣原子吸收的頻率，


\subsection{銣原子吸收譜}
為了確定相位調製對於銣原子吸收的影響，我們調整入射光的頻率，掃出整個吸收譜，如下圖黑線。
接著打開第一台 EOM，將頻寬從 30 MHz 展至 10 GHz，結果如上圖藍線，可見頻譜展寬之後，光能大部分透射銣原子氣體不被吸收。
若同時將兩台 EOM 開啟，則能再次看到光被吸收，但吸收率卻明顯降低，原因如上一小節所述，可能為訊號不夠好影響壓縮品質所致。

\section{單光子量測}
\subsection{展頻與壓縮}
以（光路圖）的架設，先不要放 $^{87}Rb$ 原子氣體管，讓單光子直接通過 60 MHz 寬的 Etalon 濾波器。若兩台 EOM 都沒開的話，窄頻的單光子能完全通濾波器，$G^{2}(\tau)$ 的量測結果如圖。此時若開啟第一台 EOM，將頻譜展至 10 GHz 寬，則量子光僅有極低的機率能通過 Etalon，如圖。若將兩台相位調製器都開啟，互補的隨機訊號能互相抵銷相位的變化，使頻譜還原至窄頻，如此就能再次通過 Etalon，如圖。

\subsection{$^{87}Rb$吸收}
同上一小節的光路架設，但把 $^{87}Rb$ 原子氣體管放回光路上。兩台相位調製器都不開的話，單光子幾乎全部被吸收，如圖。若開啟第一台相位調製器將單光子頻譜展寬，雖然光子能幾乎不被吸收，但由於 Etalon 的過濾，探測器仍測不太到光子。若將第二台相位調製器也開啟，則能把單光子的頻譜壓縮，再通過 Etalon，如圖。

單獨將圖與圖拿出來比較如下，同樣是測量展頻再壓縮回來的光，有測量

\centering 

\end{document}